\section{不足与未来展望}
\label{sec:ch07_limitations_future_work}

现有触觉观测建立在仿真模型与既定布局上，尚未完成真实传感器空间位置、物理单位及接触响应的系统标定，固定通道索引不能证明两域Taxel的物理对应。后续需要建立可追溯的标定流程，通过可控接触比较不同材料、载荷和姿态下的响应，再开展同一物体与任务条件下的仿真实机对照。可将具有明确几何来源的3D打印物体作为初步载体，逐步检验抓取候选的实机可执行性与触觉对应关系，而非仅依据观测维数判断一致性。

抓取与示教验证仍受指定物体、初态和资产版本限制，不同候选方法尚缺统一预算和判据下的公平比较，缓存与同模板变体也不能直接作为独立样本。后续应统一物理参数、计算预算和任务条件，开展跨物体、跨初态验证，并完善资产版本、Root lineage、轨迹及转换数据之间的来源追踪。视频、状态、动作和时间关联还需全流程审计，局部样本通过不能替代生产全库认证，单纯增加数据量也不足以解决质量与独立性问题。

统计伪配对存在多解与监督歧义，独立潜空间尚无认证的逐点物理语义，Natural Sim缺少同步Real目标，限制了重建评价与监督可靠性。后续应先核查统计坐标及配对前提，再研究多解的稳定性和可辨识性，在统一条件下比较单一、多候选及软监督组织方式；只有相应前提得到验证，才进一步探索潜变量结构约束与接触信息辅助监督。共享场和残差路径已实现，但仍需检验梯度优化、输入保持与生成质量的关系，不能预设增加监督形式即可改善迁移。

评价所用Real Session有限，部分Validation已有开发暴露，单Seed、容量与数值后端差异也限制了结果归因。后续需要独立Session、Root及多Seed验证，并统一模型容量、训练预算和推理后端的比较条件，结合分布、覆盖、输入依赖、跨Session稳定性及物理接触证据评价结果。当前对特定判别器的可靠性限制不应推广到所有评价器。最终仍需通过真实机器人抓取和操作成功率、接触稳定性检验触觉迁移的任务价值，现有统计指标不能替代这一验证，也不能保证未来方法获得收益。
